\section{Discussion}
\label{sec:discussion}

% Write this from the results, not from the plan. The outline below is what the
% discussion should ANSWER, not what it should conclude.

\subsection{What the presence-versus-prose result means for skill authors}
% If the SkillsBench null replicates, the practical advice inverts most current
% skill-authoring guidance: effort belongs in the description and the wiring,
% not in the body. If SkillOpt replicates instead, prose optimisation is worth
% automating and the field needs the confidence intervals it currently lacks.
% Both outcomes are actionable; say which one the data supports and how strongly.

\subsection{Where the distractor failure mode actually lives}
% If the flagship's effect is small on distractor-present items but real in
% long-horizon agentic accumulation, that relocates the failure mode rather than
% dissolving it -- and that is a more useful finding than either "it is real" or
% "it was an artifact".

\subsection{Wiring as the dominant term}
% If trigger quality moves results more than skill content does, the field's
% unit of analysis is wrong: we have been evaluating documents when we should be
% evaluating retrieval. Follow this thought honestly even though it partly
% undercuts the premise of writing skills at all.

\subsection{What a verdict is worth}
% A verdict is a claim about one model version at one point in time, under one
% harness, on tasks with computable ground truth. Given HV/MV ~ 7.8x, transfer
% to another harness should not be assumed -- which is the same warning
% arXiv:2605.23950 issues about everyone else's numbers, applied to ours.

\subsection{Negative results}
% Any NULL or HARMFUL verdict gets a written entry here: hypothesis, N, observed
% effect, CI, why we expected otherwise, and what would need to change.
%
% This is not ritual self-criticism. Publishing the nulls is what stops a
% reviewer discounting the wins, and it is the cheapest credibility available to
% a paper in a neighbourhood this crowded.
